% =====================================================================
%  7. Convergence analysis
% =====================================================================
\section{How fast does it converge?}

\begin{frame}{One inequality drives every proof}
  \begin{thmbox}[Descent lemma (proved in full in Problem M2)]
    If $\nabla f$ is $L$-Lipschitz: $\;f(\y)\le f(\x)+\nabla f(\x)\T(\y-\x)+\tfrac L2\norm{\y-\x}^2$.
  \end{thmbox}
  \small
  \textbf{Proof idea.}
  \pause
  $f(\y)-f(\x)-\nabla f(\x)\T(\y-\x) = \int_0^1[\nabla f(\x+t(\y-\x))-\nabla f(\x)]\T(\y-\x)\,dt$,
  \pause
  and Cauchy--Schwarz with the Lipschitz bound turns this into $\le\int_0^1 Lt\norm{\y-\x}^2dt = \tfrac L2\norm{\y-\x}^2$. $\blacksquare$

  \pause
  \vspace{4pt}
  \normalsize Plug in $\y = \x_k - \alpha\g_k$:
  \[
    f(\x_{k+1}) \le f(\x_k) - \alpha\Bigl(1 - \frac{\alpha L}2\Bigr)\norm{\g_k}^2
    \;\overset{\alpha = 1/L}{=}\; f(\x_k) - \frac{1}{2L}\norm{\g_k}^2 .
  \]
  \small Each step buys a guaranteed drop of $\norm{\g_k}^2/(2L)$. Every rate that follows comes from linking
  $\norm{\g_k}$ to how far we still are from the optimum.
\end{frame}

\begin{frame}{No convexity: we can only promise a small gradient}
  \begin{thmbox}[Theorem ($L$-smooth, bounded below, $\alpha = 1/L$)]
    \[
      \min_{0\le k<K}\norm{\nabla f(\x_k)} \le \sqrt{\frac{2L\,(f(\x_0)-f^\star)}{K}} = O\!\bigl(1/\sqrt K\bigr).
    \]
  \end{thmbox}
  \small\textbf{Proof.} Add up the per-step drops for $k = 0,\dots,K-1$. The total cannot exceed
  $f(\x_0) - f(\x_K)\le f(\x_0)-f^\star$ (details in Problem M2). $\blacksquare$
  \begin{itemize}\small
    \item All we get is a point with a \emph{small gradient}. It could be a saddle or a poor local minimum.
    \item Reaching $\norm{\nabla f}\le\varepsilon$ takes $O(L(f_0-f^\star)/\varepsilon^2)$ iterations, which is slow.
          That is what you pay for assuming nothing.
    \item Deep learning lives here. The faster rates on the next slides all need extra structure.
  \end{itemize}
\end{frame}

\begin{frame}{Convex functions: error $O(1/k)$}
  \begin{thmbox}[Theorem ($L$-smooth, convex, $\alpha = 1/L$)]
    \[ f(\x_k) - f^\star \le \frac{L\norm{\x_0-\x^\star}^2}{2k}. \]
  \end{thmbox}
  \small\textbf{Proof.}
  \pause
  Convexity gives $f(\x_k)\le f^\star + \g_k\T(\x_k-\x^\star)$. Combine that with the guaranteed drop:
  \[
    f(\x_{k+1}) - f^\star \le \g_k\T(\x_k-\x^\star) - \tfrac{1}{2L}\norm{\g_k}^2
    = \tfrac L2\Bigl(\norm{\x_k-\x^\star}^2 - \norm{\x_{k+1}-\x^\star}^2\Bigr).
  \]
  \pause
  (The equality is just $\norm{\x_k - \x^\star - \g_k/L}^2$ expanded.) Summing over $k = 0,\dots,K-1$, the right side
  telescopes:
  \[
    \textstyle\sum_{k=1}^{K}\bigl(f(\x_k)-f^\star\bigr)\le\tfrac L2\norm{\x_0-\x^\star}^2 .
  \]
  \pause
  The values $f(\x_k)$ never increase, so the left side is at least $K\bigl(f(\x_K)-f^\star\bigr)$. $\blacksquare$
\end{frame}

\begin{frame}{Strongly convex functions: the error shrinks geometrically}
  \begin{thmbox}[Theorem ($L$-smooth, $\mu$-strongly convex, $\alpha = 1/L$)]
    \[ \norm{\x_{k}-\x^\star}^2 \le \Bigl(1-\frac{\mu}{L}\Bigr)^{k}\norm{\x_0-\x^\star}^2 = \Bigl(1-\frac1\kappa\Bigr)^k\norm{\x_0-\x^\star}^2. \]
  \end{thmbox}
  \small\textbf{Proof.} Write $\e = \x_k-\x^\star$, $\Delta = f(\x_k)-f^\star$, and expand
  $\norm{\x_{k+1}-\x^\star}^2 = \norm{\e}^2 - \tfrac2L\g_k\T\e + \tfrac1{L^2}\norm{\g_k}^2$. Two facts bound the middle and last terms:
  \pause
  \begin{itemize}\footnotesize
    \item strong convexity at $\y=\x^\star$: $\;\g_k\T\e\ge\Delta+\tfrac\mu2\norm{\e}^2$;
    \item the descent lemma at $\x_k-\g_k/L$: $\;f^\star\le f(\x_k)-\tfrac1{2L}\norm{\g_k}^2$, i.e. $\norm{\g_k}^2\le2L\Delta$.
  \end{itemize}
  \pause
  \[
    \norm{\x_{k+1}-\x^\star}^2\le\norm{\e}^2-\tfrac2L\Delta-\tfrac\mu L\norm{\e}^2+\tfrac2L\Delta
    =\Bigl(1-\tfrac\mu L\Bigr)\norm{\e}^2. \quad\blacksquare
  \]
  \pause
  \footnotesize With $\alpha = \frac{2}{\mu+L}$ the factor improves to $\norm{\e_{k+1}}\le\frac{\kappa-1}{\kappa+1}\norm{\e_k}$,
  the same number the quadratic analysis gave in Section~5~\cite{nesterov2018}.
\end{frame}

\begin{frame}{Linear rates without convexity: the Polyak--{\L}ojasiewicz condition}
  \begin{thmbox}[PL inequality \normalfont\cite{karimi2016}]
    $\tfrac12\norm{\nabla f(\x)}^2\ge\mu\,(f(\x)-f^\star)$ for all $\x$. Then GD with $\alpha = 1/L$ satisfies
    \[ f(\x_k)-f^\star\le\Bigl(1-\frac\mu L\Bigr)^k\bigl(f(\x_0)-f^\star\bigr). \]
  \end{thmbox}
  \small\textbf{Proof (two lines).}
  $f(\x_{k+1})-f^\star\le f(\x_k)-f^\star-\tfrac1{2L}\norm{\g_k}^2\le\bigl(1-\tfrac\mu L\bigr)(f(\x_k)-f^\star)$. $\blacksquare$
  \begin{itemize}\small
    \item PL does \emph{not} need convexity. $f(x) = x^2+3\sin^2x$ has $f''$ dipping to $-4$,
          yet it satisfies PL (numerically $\mu\approx0.175$).
    \item Strong convexity implies PL (from the Section~6 sandwich), so PL is the more general route to a linear rate.
    \item Least squares $\tfrac12\norm{\mathbf M\x-\vb}^2$ with a rank-deficient $\mathbf M$ is convex but not strongly convex.
          It still satisfies PL, so GD still converges linearly in $f$~\cite{karimi2016}.
  \end{itemize}
\end{frame}

\begin{frame}{Order of convergence, and the scorecard so far}
  \begin{columns}[T]
    \begin{column}{0.55\textwidth}
      \small
      A sequence $\{\e_k\}$ has \textbf{order $p$} if $\norm{\e_{k+1}}\approx r\norm{\e_k}^p$.
      \begin{itemize}\footnotesize
        \item GD has $p = 1$, \alert{linear} order, with $r\le\rho$. It gains the \emph{same} number of correct digits,
              $-\log_{10}\rho$, every step.
        \item In the plot, $\log e_{k+1}$ against $\log e_k$ is a line of slope 1. The dashed slope-2 line
              (quadratic order, as in Newton-type methods) is there for comparison.
      \end{itemize}
      \vspace{2pt}
      \footnotesize
      \begin{tabular}{@{}l l l@{}}
        \toprule
        Setting & Guarantee ($\alpha = \frac1L$) & Iterations \\
        \midrule
        nonconvex & $\min\norm{\g}^2\le\frac{2L\Delta_0}{K}$ & $O(\varepsilon^{-2})$ \\
        convex & $f-f^\star\le\frac{LR^2}{2k}$ & $O(\varepsilon^{-1})$ \\
        strongly convex & $\norm{\e_k}^2\le(1-\frac1\kappa)^k R^2$ & $O(\kappa\log\frac1\varepsilon)$ \\
        PL & $f-f^\star\le(1-\frac1\kappa)^k\Delta_0$ & $O(\kappa\log\frac1\varepsilon)$ \\
        \bottomrule
      \end{tabular}
    \end{column}
    \begin{column}{0.43\textwidth}
      \centering
      \includegraphics[width=\linewidth]{fig_order}

      \footnotesize\textcolor{mGrey}{$R = \norm{\x_0-\x^\star}$ and $\Delta_0 = f(\x_0)-f^\star$.}
    \end{column}
  \end{columns}
\end{frame}

\begin{frame}{GD or Newton? Count the cost}
  \small
  \emph{When GD, and when Newton?} Count the work for $n$ parameters~\cite{nocedal2006}.

  \vspace{2pt}
  \centering\footnotesize
  \begin{tabular}{@{}l l l@{}}
    \toprule
    & Gradient descent & Newton's method \\
    \midrule
    information used & gradient ($n$ numbers) & gradient + Hessian ($n^2$ numbers) \\
    cost of one step & $O(n)$ after the gradient & $O(n^3)$ to solve $H\mathbf y = \g$ \\
    order of convergence & linear & quadratic \\
    learning rate & must be chosen & not needed \\
    on a concave stretch ($f'' < 0$) & still goes downhill & can head for a \textcolor{mRed}{maximum} \\
    best for & large problems (ML, $n\sim10^6$) & small problems needing high accuracy \\
    \bottomrule
  \end{tabular}

  \raggedright
  \vspace{3pt}
  \begin{keybox}[Why ML uses GD]
    \footnotesize With $n = 10^6$ a Hessian has $10^{12}$ entries and will not fit in memory; GD stores $n$ numbers.
    When $H$ is not positive definite, Levenberg--Marquardt~\cite{nocedal2006} moves towards a gradient step.
  \end{keybox}
\end{frame}
